\section{Conclusion and Outlook}
\label{sec:conclusion}

A general-purpose foundation embedding loses to a dataset-fitted linear reduction
here, most plausibly because the label is defined per cell line and the frozen embedding carries less
between-line structure on this atlas, and with
\NLinesCV{} labels a ridge regression on cell-line means matches everything built in this work: the
ceiling is the data, and the analysis identifies which part of the data sets it.

The next steps address the label bottleneck, the model's objective and the third research question, in
that order of directness.

\paragraph{MIL as the primary model, with a ranking loss.} A model trained on bags of a line's cells
leaves each cell free to differ (Equation~\ref{eq:mil-decomposition}), and since the metric ranks lines
within a compound, a ranking loss would align the objective with what is scored
\cite{ilse2018,burges2005}.

\paragraph{Bulk pretraining and domain adaptation.} CTRPv2 alone covers \NCtrpLines{} cell lines
against the \NLines{} here, so labels exist and single-cell expression is what is missing
\cite{ctrpv2,scdeal}. Pretraining on bulk expression and aligning bulk and single-cell spaces by domain
adaptation addresses that gap \cite{scdeal,ganin2015}. The most direct form is to apply PERCEPTION's
bulk-pretrained response models to these cells and compare their predictions with those from the two
representations here \cite{perception}.

\paragraph{Fine-tuning scGPT, given more data.} Fine-tuning tests whether the ordering against PCA
reverses once the encoder may keep line identity, but on about \NLinesCV{} labels a fine-tuned encoder
may memorise them, so it follows the previous step.

\paragraph{Explicit explainability of heterogeneity.} The question could be reframed from implicit
learning to explicit explanation: identifying the subpopulations within cell lines that drive resistance
and relating them to the recurrent heterogeneity programs annotated for this atlas \cite{scp542}.

\paragraph{Missing controls.} Every comparison here is between two reductions of one matrix. A
line-level model on the \NHvg{} selected genes alone would show whether either reduction helps at all. Division rate is a further bias to extract, as fragility is: a 72-hour viability assay
confounds compound effect with proliferation \cite{hafner}, and an earlier test against the dataset's
cell-cycle scores, negative but run on a since-retired target, used a snapshot rather than a rate, so
measured doubling times are the open version. The input-dropout rate could be matched in effect rather
than in value, by choosing per arm the rate that removes the same expected variance
\cite{srivastava2014}.
